\documentclass[11pt,twoside]{article}

\usepackage[margin=1in]{geometry}
\usepackage{amsmath,amssymb,amsthm,mathtools}
\usepackage[inline,shortlabels]{enumitem}
\usepackage{booktabs}
\usepackage{xcolor}
\usepackage{microtype}

% Revision date: maintained by scripts/build_lecture_notes.py. Keep this command.
\newcommand{\lecturerevised}{2026-09-28 20:44 UTC-06:00}
\usepackage{eso-pic}
\AddToShipoutPictureBG{%
  \AtPageLowerLeft{%
    \put(\LenToUnit{0.5\paperwidth},\LenToUnit{0.5in}){%
      \makebox(0,0){\normalfont\footnotesize\color{black!60}Last revised: \lecturerevised}%
    }%
  }%
}

\usepackage[square,authoryear]{natbib}
\usepackage[colorlinks=true,linkcolor=blue!60!black,
  urlcolor=blue!60!black,citecolor=blue!60!black]{hyperref}
\usepackage[nameinlink,capitalize]{cleveref}
\crefname{exercise}{Exercise}{Exercises}
\Crefname{exercise}{Exercise}{Exercises}
\hypersetup{pdftitle={Lecture 7: Distribution shift and helping models learn},
  pdfauthor={Csaba Szepesvari}}

\usepackage{../exercises/course-exercises}

\usepackage{enotez}
\DeclareInstance{enotez-list}{lecture}{list}{format=\normalsize,number={#1.}}
\setenotez{list-name=Endnotes,list-style=lecture,backref=true}
\crefname{endnote}{Endnote}{Endnotes}
\Crefname{endnote}{Endnote}{Endnotes}

\definecolor{courseblue}{HTML}{17365D}
\newcommand{\E}{\mathbb E}
\newcommand{\I}{\mathbb I}
\renewcommand{\P}{\mathbb P}
\newcommand{\Law}{\mathcal L}
\newcommand{\R}{\mathbb R}
\newcommand{\KL}{D}
\newcommand{\TV}{\mathrm{TV}}
\makeatletter
\@for\@letter:=A,B,C,D,E,F,G,H,I,J,K,L,M,N,O,P,Q,R,S,T,U,V,W,X,Y,Z\do{%
  \expandafter\edef\csname c\@letter\endcsname{\noexpand\mathcal{\@letter}}}
\makeatother

\newtheorem{theorem}{Theorem}[section]
\newtheorem{proposition}[theorem]{Proposition}
\newtheorem{corollary}[theorem]{Corollary}
\newtheorem{lemma}[theorem]{Lemma}
\theoremstyle{definition}
\newtheorem{definition}[theorem]{Definition}
\newtheorem{example}[theorem]{Example}
\theoremstyle{remark}
\newtheorem{remark}[theorem]{Remark}

\renewcommand{\thepage}{7--\arabic{page}}
\renewcommand{\thesection}{7.\arabic{section}}
\renewcommand{\theequation}{7.\arabic{equation}}
\pagestyle{myheadings}
\markboth{Lecture 7: Distribution shift and helping models learn}
{Lecture 7: Distribution shift and helping models learn}

\begin{document}

\begin{center}
\fcolorbox{courseblue}{white}{%
\begin{minipage}{0.92\textwidth}
\vspace{1mm}
\textbf{CMPUT 654: Mathematical Foundations of Modern AI Systems}
\hfill Fall 2026

\vspace{3mm}
\begin{center}
{\Large\bfseries Lecture 7: Distribution shift and helping models learn}

\vspace{1mm}
September 24, 2026
\end{center}

\vspace{2mm}
\textit{Lecturer: Csaba Szepesv\'ari}
\hfill
\textit{Note prepared by: Csaba Szepesv\'ari}
\vspace{1mm}
\end{minipage}}
\end{center}

\noindent\textbf{Reading guide.}
The guiding question is: when does a change in the data law change what should
be learned, and when can a change in the supplied data make learning easier?
The first part distinguishes covariate shift from general distribution shift
and explains why correct specification and coverage determine whether a shift
is harmful. The second part organizes data selection as communication between
a teacher and a student. The third part uses parity to illustrate a difficulty
in teaching a neural network a desired behavior. Direct final-answer training
can fail even when the network can represent the target. A modified training
task exposes intermediate computations while preserving the desired final
answer. The lecture stated the negative result, explained its proof mechanism,
and previewed the intermediate-supervision construction. Lecture 8 supplies
the complete gradient and trajectory argument and develops the positive
constructions.
The prerequisites are conditional expectation, KL divergence, total
variation, Pinsker's inequality, Hoeffding's inequality, ReLU networks, and
elementary linear algebra. The brief finite-field and Fourier explanations
below supply the additional background used in the parity proof.

\section{Distribution shift depends on the prediction problem}
\label{sec:shift}

Let $P_{\rm tr}$ and $P_{\rm dep}$ be the training and deployment laws on
$\cX\times\cY$. Write $\mu_{\rm tr}$ and $\mu_{\rm dep}$ for their input
marginals.

For a fixed predictor $f$, let $c_f:\cX\times\cY\to[0,1]$ be its downstream
cost and write
\[
 J_P(f)=\E_P[c_f(X,Y)].
\]

\begin{proposition}[Total variation controls every fixed bounded cost]
\label{prop:tv-fixed-cost}
For every fixed predictor $f$,
\begin{equation}
 J_{P_{\rm dep}}(f)-J_{P_{\rm tr}}(f)
 \le \TV(P_{\rm dep},P_{\rm tr}).
\label{eq:tv-fixed-cost-one-sided}
\end{equation}
Consequently,
\begin{equation}
 \left|J_{P_{\rm dep}}(f)-J_{P_{\rm tr}}(f)\right|
 \le \TV(P_{\rm dep},P_{\rm tr}).
\label{eq:tv-fixed-cost}
\end{equation}
\end{proposition}

\begin{proof}
Take a measure $\nu$ dominating both laws and let $p_{\rm dep}$ and
$p_{\rm tr}$ be their densities. Then
\begin{align*}
J_{P_{\rm dep}}(f)-J_{P_{\rm tr}}(f)
&=\int c_f(z)\bigl(p_{\rm dep}(z)-p_{\rm tr}(z)\bigr)\,\nu(dz)\\
&\le \int_{\{p_{\rm dep}\ge p_{\rm tr}\}}
       \bigl(p_{\rm dep}(z)-p_{\rm tr}(z)\bigr)\,\nu(dz)\\
&=\TV(P_{\rm dep},P_{\rm tr}).
\end{align*}
Interchanging the two laws gives the absolute-value bound.
\end{proof}

The bound is uniform over all costs in $[0,1]$. For the particular cost of a
particular predictor, the exact difference is the signed integral in the first
line of the proof. It can be small even when the total-variation distance is
large, because the cost can be small where the laws differ.

\begin{definition}[Covariate shift and general distribution shift]
There is \emph{covariate shift} when the two joint laws admit the
factorizations
\begin{equation}
 P_{\rm tr}(dx,dy)=\mu_{\rm tr}(dx)K(dy\mid x),
 \qquad
 P_{\rm dep}(dx,dy)=\mu_{\rm dep}(dx)K(dy\mid x)
 \label{eq:covariate-shift}
\end{equation}
for one conditional kernel $K$, while $\mu_{\rm tr}\ne\mu_{\rm dep}$.
This definition also specifies the conditional target at inputs that occur
under only one of the two marginals. There is \emph{general distribution shift}
when $P_{\rm tr}\ne P_{\rm dep}$; the conditional law of $Y$ given $X$ may
change as well.
\end{definition}

Covariate shift changes how often inputs occur while preserving the desired
prediction at each input. General distribution shift may also change the
desired prediction. This classification depends on what is included in $X$.

\begin{proposition}[A source indicator can turn conditional shift into covariate shift]
Suppose that $X,Y$, and the source indicator $J$ take values in finite sets.
For each source $j$, let $\mu_j$ be its input law and $K_j$ its conditional
kernel, so that $P_j(x,y)=\mu_j(x)K_j(y\mid x)$. For mixture weights
$\alpha_j$, define the joint law
\[
 P_\alpha(x,y,j)=\alpha_jP_j(x,y),
 \qquad \alpha_j\ge0,
 \qquad \sum_j\alpha_j=1.
\]
If the predictor observes only $X$, its conditional kernel is, at every $x$
with positive $P_\alpha$-probability,
\begin{equation}
 K_\alpha(y\mid x)=\sum_jK_j(y\mid x)P_\alpha(j\mid x),
 \label{eq:source-hidden}
\end{equation}
which generally changes with the mixture weights. If the predictor observes
$(X,J)$, then
\begin{equation}
 \widetilde K_\alpha(y\mid x,j)=K_j(y\mid x),
 \label{eq:source-visible}
\end{equation}
which is independent of the mixture weights.
\end{proposition}

The same corpus change can therefore be general shift for one prediction
problem and covariate shift for another. Source metadata is useful only if the
model receives it, preserves it, or can infer it from the context.

\begin{example}[Benign and harmful changes]
Increasing the proportion of mathematics prompts is covariate shift if the
desired answer to every prompt is unchanged. Combining sources that use
incompatible conventions changes the conditional law of $Y$ given $X$ when
the source is hidden. A
recommender system can create a further effect: its outputs change what users
consume, which can later change their preferences and hence the conditional
law of feedback.
\end{example}

Evaluation has the same dependence on its law. For a loss $\ell$, an
evaluation set drawn from $P_{\rm eval}$ estimates
\[
 L_{P_{\rm eval}}(f)=\E_{P_{\rm eval}}[\ell(f(X),Y)].
\]
It estimates deployment risk when its construction represents the deployment
law, or when a separate argument relates the two laws. Curation is therefore
part of evaluation as well as training.

\section{Correct specification, misspecification, and coverage}
\label{sec:specification}

For a joint law $P$ on $\cX\times\cY$, let $X,Y$ denote its coordinate maps
and define the population risk
\[
 L_P(f)=\E_P[\ell(f(X),Y)].
\]

\begin{definition}[Correct specification and misspecification]
A conditional model class $\cF$ is \emph{correctly specified} for $P$ if it
contains a predictor that attains the optimal conditional decision at every
input $x\in\cX$. It is \emph{misspecified} if no member of the class does so.
\end{definition}

Correct specification is a claim about the interaction among the target
conditional law, the loss, and the model class. It is stronger than saying
that the class contains a predictor with small average training risk.

\begin{proposition}[A common pointwise optimum survives covariate shift]
\label{prop:pointwise-optimum}
Suppose $P$ and $Q$ share a conditional kernel $K(dy\mid x)$ as in
\cref{eq:covariate-shift}, and write $\mu_P$ and $\mu_Q$ for their input
marginals. If $f^\star\in\cF$ satisfies, for every
$x\in\cX$ and every $f\in\cF$,
\[
 \int \ell(f^\star(x),y)K(dy\mid x)
 \le \int \ell(f(x),y)K(dy\mid x),
\]
then $f^\star$ minimizes both $L_P$ and $L_Q$.
\end{proposition}

\begin{proof}
\begin{align*}
L_Q(f)-L_Q(f^\star)
&=\int_{\cX}\left[
  \int_{\cY}\bigl(\ell(f(x),y)-\ell(f^\star(x),y)\bigr)K(dy\mid x)
  \right]\mu_Q(dx)\\
&\ge0.
\end{align*}
Replacing $\mu_Q$ by $\mu_P$ proves the same claim for $P$.
\end{proof}

For conditional log loss, let $p^\star(\cdot\mid x)$ be the true conditional
law and $q_f(\cdot\mid x)$ the model prediction under $P$.
\begin{equation}
L_P(f)=H_P(Y\mid X)+
\E_P\!\left[
\KL\bigl(p^\star(\cdot\mid X)\Vert q_f(\cdot\mid X)\bigr)
\right].
\label{eq:conditional-logloss}
\end{equation}
If one model represents $p^\star(\cdot\mid x)$ at every deployment input,
changing the input mixture does not change this population optimum. The
mixture can still change the number of observations at each input, the
conditioning of the optimization problem, and which optimum is selected in an
overparameterized class.

\begin{proposition}[Misspecification makes the mixture part of the target]
\label{prop:misspecified-mixture}
Let $\cX=\{a,b\}$, $\cY=\{0,1\}$, and define a conditional kernel $K$ by
\[
K(1\mid a)=0.1,
\qquad
K(1\mid b)=0.9.
\]
For $\lambda\in[0,1]$, let $\mu_\lambda(a)=\lambda$,
$\mu_\lambda(b)=1-\lambda$, and
$P_\lambda(x,y)=\mu_\lambda(x)K(y\mid x)$. Restrict the model to one
probability $q$ at both inputs. The log-loss minimizer under $P_\lambda$ is
\[
q_\lambda^\star=0.9-0.8\lambda.
\]
\end{proposition}

\begin{proof}
Write $r_\lambda=P_\lambda(Y=1)=0.9-0.8\lambda$. The risk is
\[
L_{P_\lambda}(q)
=-r_\lambda\log q-(1-r_\lambda)\log(1-q).
\]
Therefore
\[
\frac{d}{dq}L_{P_\lambda}(q)
=-\frac{r_\lambda}{q}+\frac{1-r_\lambda}{1-q},
\]
and the unique stationary point is $q=r_\lambda$.
\end{proof}

The restricted model must compromise between the two contexts. Their weights
determine the compromise. In a sufficiently expressive language model, the
same parameters may represent both mathematics and literature conditionals,
so a moderate mixture change need not change the ideal predictor. This is an
inductive-bias claim about the shared representation; it does not follow from
model size alone.

A \emph{minimum-coverage condition} supplies a separate requirement. The
density-ratio condition below is one sufficient, model-independent form of
coverage; it is stronger than necessary for structured model classes.

\begin{proposition}[Minimum coverage transfers excess risk]
\label{prop:coverage}
Assume covariate shift with input marginals $\mu_{\rm tr}$ and
$\mu_{\rm dep}$, and let
\[
\Delta_f(x)=
\int\bigl(\ell(f(x),y)-\ell(f^\star(x),y)\bigr)K(dy\mid x)\ge0.
\]
If $\rho\in(0,1]$, $\mu_{\rm dep}$ is absolutely continuous with respect to
$\mu_{\rm tr}$ and
\[
\frac{d\mu_{\rm dep}}{d\mu_{\rm tr}}(x)\le\rho^{-1}
\quad \mu_{\rm tr}\text{-almost surely},
\]
then
\begin{equation}
L_{P_{\rm dep}}(f)-L_{P_{\rm dep}}(f^\star)
\le
\rho^{-1}\!\left(
L_{P_{\rm tr}}(f)-L_{P_{\rm tr}}(f^\star)
\right).
\label{eq:coverage-transfer}
\end{equation}
\end{proposition}

\begin{proof}
Let $\omega=d\mu_{\rm dep}/d\mu_{\rm tr}$. Then
\[
\int\Delta_f(x)\,\mu_{\rm dep}(dx)
=\E_{\mu_{\rm tr}}[\omega\Delta_f]
\le\rho^{-1}\E_{\mu_{\rm tr}}[\Delta_f].
\]
\end{proof}

If deployment assigns positive mass to a region absent from training, no
finite $\rho$ exists. This does not make correct prediction there impossible.
A structured class can identify behavior at unseen inputs from observations
elsewhere. For example, a linear predictor can be identified when the training
covariates cover the parameter directions needed at deployment, even if the
individual deployment inputs were never observed. Thus the necessary notion
of coverage depends on the model and the prediction problem: it must identify
the relevant degrees of freedom, rather than literally include every
deployment input. The density-ratio condition avoids this structural analysis
by imposing a stronger distribution-level requirement.
The term \emph{domain adaptation} commonly refers to methods that use source
data together with information about a different target domain, often
unlabeled target inputs, to obtain a predictor for the target law.

\section{Learning as communication between teacher and student}
\label{sec:teacher-student}

Data selection can be described as a communication protocol. A
\emph{teacher} sends data or feedback; a \emph{student} updates its internal
state and may send queries or predictions. A useful classification separates
four questions.

\begin{center}
\begin{tabular}{p{0.22\textwidth}p{0.68\textwidth}}
\toprule
Dimension & Question \\
\midrule
Objective information & Does the teacher know the target predictor, desired
behavior, or only a verifier? What does the student know? \\
Student information & Does the teacher know the student's representation or
hypothesis class, update algorithm, current state, or only observed outputs? \\
Control of communication & Who chooses the next input, label, intermediate
state, correction, or evaluation: the teacher, the student, or both? \\
Number of turns & Is the data fixed in one batch, or can later messages depend
on the transcript and the student's current state? \\
\bottomrule
\end{tabular}
\end{center}

These dimensions are independent. In a \emph{passive learning} protocol, the
student receives data selected without using its queries or current state.
\emph{Data curation} selects examples or changes their weights;
\emph{task curation} changes the supervised task; and a \emph{curriculum}
controls the order or difficulty of the supplied examples or tasks. Any of
these choices can be fixed in advance or adapted over several turns.

\emph{Active learning} usually gives the student control over queries whose
answers are supplied by an oracle. In \emph{active teaching}, the teacher
chooses the next message using the interaction so far. \emph{Machine teaching}
usually gives the teacher a target and a model of the student, and asks the
teacher to choose effective training data. In a multi-turn protocol, either
side or both sides can adapt to the transcript. Thus ``active'' identifies who
selects the next message in the protocol.

The teacher's information is a resource. Knowing the target and the student's
update rule can reduce the required number of examples. It can also make a
teaching strategy brittle when the assumed student differs from the actual
one. The student's ability to query is another resource. It can expose the
specific uncertainty that the current data leave unresolved.

The communication channel also matters. Reordering final-answer examples
changes the optimization path. Replacing
\[
(x,y)
\quad\text{by}\quad
(x,z_1,\ldots,z_r,y)
\]
adds information about intermediate computations. Mathematical reasoning is
a favorable setting because problems can often be generated, final answers
can often be verified, intermediate states may be checkable, and computation
length can be controlled.

When learning is modeled through a hypothesis class, correct specification
says that the desired behavior is represented. It does not imply that the
student's update rule can find that behavior. Changing the communicated data
or the supervised task can change this optimization problem while preserving
the desired final output.

The parity example applies this distinction between the desired behavior and
the training task used to teach it. The desired final output remains parity.
What the teacher reveals during training will change.

\clearpage
\section{Parity as a test case for teaching a computation}
\label{sec:parity-setup}

This example illustrates one difficulty in teaching a neural network a
desired behavior and one way to overcome it. The direct approach supplies
input--output pairs and trains the network to predict only the final answer.
The result in \cref{sec:parity-result} shows that a specified version of this
approach can fail. This does not imply that every way of teaching a neural
network to produce the same answer must fail. In \cref{sec:intermediate-preview}, the
training task will expose intermediate computations, and completing the
modified task will still produce the original final answer.

The same parity problem has two useful representations. For
$b\in\{0,1\}^d$, define
\[
 \operatorname{parity}_S(b)=\bigoplus_{j\in S}b_j.
\]
Set $x_j=(-1)^{b_j}$. Then
\begin{equation}
 (-1)^{\operatorname{parity}_S(b)}
 =\prod_{j\in S}x_j.
\label{eq:parity-representations}
\end{equation}
The bit representation makes XOR computation and linear algebra over
$\mathbb F_2$ explicit. The sign representation makes products, correlations,
and the ReLU gradient calculation explicit.

For $S\subseteq[d]$ and $x\in\{-1,1\}^d$, define
\[
\chi_S(x)=\prod_{j\in S}x_j.
\]
Fix $S$ and define the data law $P_S$ on
$\{-1,1\}^d\times\{-1,1\}$ by
\[
 P_S(x,y)=2^{-d}\I\{y=\chi_S(x)\}.
\]
Let $X,Y$ be the coordinate maps under $P_S$. Thus $X$ is uniform on the
Boolean cube and $Y=\chi_S(X)$ with $P_S$-probability one.

We note in passing that for any $0<\delta<1$,
after seeing $\Omega(d+ \log(1/\delta))$ labeled examples
where the inputs come from the uniform distribution,
$S$ can be determined with probability $1-\delta$ in a computationally
efficient manner (see
\cref{ex:l7-parity-identification}).
As we shall see, training a neural network with gradient descent
to do the same does not succeed. However, before that, we establish that the target
lies  in the network class considered in the negative result: without this result,
the negative result would not make sense.

\subsection{A small ReLU network represents the fixed parity}

Consider the one-hidden-layer ReLU network
\begin{equation}
N_\theta(x)=\sum_{r=1}^m u_r[w_r^\top x+b_r]_+.
\label{eq:relu-network}
\end{equation}

\begin{proposition}[Parity has a small ReLU representation]
\label{prop:parity-representation}
For every $S$, a network of the form \cref{eq:relu-network} with width at
most $|S|+1$ represents $\chi_S$ on the Boolean cube.
\end{proposition}

\Cref{ex:l7-parity-representation} asks you to verify an explicit
width-$(|S|+1)$ construction.

\section{Direct gradient training can fail for a fixed parity}
\label{sec:parity-result}

For a score $f$, define its \emph{signed correlation loss}
\[
\cC(f)=-\E_{P_S}[Yf(X)],
\]
and define the parameter objective
\begin{equation}
F(\theta)=\cC(N_\theta)=-\E_{P_S}[Y N_\theta(X)].
\label{eq:correlation-objective}
\end{equation}
For a hard predictor $h:\{-1,1\}^d\to\{-1,1\}$, define its classification
risk by $R_{01}(h)=P_S(h(X)\ne Y)$. For binary-valued predictions,
$\cC(h)=2R_{01}(h)-1$. The identity does not extend to a real-valued score
$f$: an arbitrarily small score can have the correct sign everywhere.
\Cref{ex:l7-correlation-loss} asks you to prove these claims and derive the
corresponding decompositions for squared and logistic loss. The main result
below analyzes perturbed gradient descent on $F$ itself. Its classification
conclusion follows from symmetry of the entire parameter law, rather than
from the numerical value of $\cC(N_\theta)$.

On an ambient probability space, let $\Theta_0,\Xi_0,\ldots,\Xi_{T-1}$ be
mutually independent, with
\[
 \Law(\Theta_0)=\cN(0,\sigma^2I),
 \qquad
 \Law(\Xi_t)=\cN(0,\sigma^2I).
\]
Starting from $\Theta_0$, run \emph{perturbed gradient descent}
\begin{equation}
\Theta_{t+1}=\Theta_t-\eta\nabla F(\Theta_t)-\Xi_t.
\label{eq:perturbed-gd}
\end{equation}
Define the comparison process obtained by deleting the gradient:
\begin{equation}
\Theta^{(0)}_0=\Theta_0,
\qquad
\Theta^{(0)}_{t+1}=\Theta^{(0)}_t-\Xi_t.
\label{eq:gaussian-walk}
\end{equation}
The independent Gaussian increments make the comparison process Gaussian at
every time. This symmetry is the reason the small-gradient estimate can be
applied repeatedly along the comparison path.

The proof of the main result uses the following comparison between the trained
parameter path and the Gaussian walk.

\begin{lemma}[Path comparison]
\label{lem:path-comparison}
Let $d>30$, $k=|S|$, and $A>0$. Let
\(m,T\in\mathbb N\) and \(\eta,\sigma>0\), and assume
\begin{equation}
m,T,\eta,\sigma,\sigma^{-1}\le d^A,
\qquad
k\ge(288A+144)\log d.
\label{eq:hardness-conditions}
\end{equation}
Then
\begin{equation}
\TV\!\left(
\Law(\Theta_0,\ldots,\Theta_T),
\Law(\Theta^{(0)}_0,\ldots,\Theta^{(0)}_T)
\right)
\le e^{-k/36}.
\label{eq:path-tv}
\end{equation}
\end{lemma}

The comparison process is independent of $S$. The lemma says that the
gradient changes the law of a polynomial-length parameter path by only an
exponentially small amount, even though the network class contains the target.

The symmetry of the comparison process now gives the main result.

\begin{theorem}[Classification advantage occurs with probability at most one half]
\label{thm:classification}
Let $d>30$, $k=|S|$, and $A>0$. Let
\(m,T\in\mathbb N\) and \(\eta,\sigma>0\), and assume
\[
m,T,\eta,\sigma,\sigma^{-1}\le d^A,
\qquad
k\ge(288A+144)\log d.
\]
Let $h_\theta(x)=\operatorname{sign}(N_\theta(x))$, with fair tie-breaking
when $N_\theta(x)=0$. Equivalently,
\[
R_{01}(h_\theta)
=P_S(YN_\theta(X)<0)+\frac12P_S(N_\theta(X)=0).
\]
Then, for every $\gamma>0$,
\begin{equation}
\P\!\left(R_{01}(h_{\Theta_T})\le\frac12-\gamma\right)
\le\frac12+e^{-k/36},
\label{eq:classification-probability}
\end{equation}
and
\begin{equation}
\left|\E[R_{01}(h_{\Theta_T})]-\frac12\right|
\le e^{-k/36}.
\label{eq:classification-expectation}
\end{equation}
\end{theorem}

The expectation bound permits a symmetric mixture of very accurate and very
inaccurate classifiers. The probability bound gives the corresponding
reliability statement: the chance of achieving any fixed positive advantage
over random guessing is at most one half, up to the exponentially small term.


\section{Why the comparison theorem is plausible}
\label{sec:proof-roadmap}

The proof of \cref{lem:path-comparison} has two parts. First, when the
parameters have an isotropic Gaussian law, the expected norm of the
parity-dependent gradient is exponentially small in $k$. The calculation
expresses each gradient coordinate through correlations between a random
threshold and a degree-$k$ Boolean Fourier character. Low-degree terms cancel,
leaving only an exponentially small remainder.

Second, the comparison process in \cref{eq:gaussian-walk} is Gaussian at every
time. The small-gradient estimate can therefore be applied repeatedly along
this path. A KL chain-rule calculation compares the actual update kernel with
an auxiliary kernel that deletes sufficiently small gradients. Pinsker's
inequality and a coupling of the auxiliary process with the Gaussian walk then
give \cref{eq:path-tv}. The sign symmetry of the Gaussian walk yields the
classification statement in \cref{thm:classification}.

This argument explains the obstruction. The network class contains the fixed
parity, but typical Gaussian parameters lie in a region where the objective
supplies exponentially little information about the direction in which the
parameters should move. Lecture 8 gives the complete calculation.

\section{Changing the training task}
\label{sec:intermediate-preview}

The lower bound concerns direct training on final parity labels. It does not
apply to every way of teaching a neural network to compute parity. Write
$x_i=(-1)^{b_i}$ with $b_i\in\{0,1\}$ and define the running parities
\[
z_1=b_1,
\qquad
z_i=z_{i-1}\oplus b_i.
\]
The final value $z_d$ determines the original parity, while every intermediate
label demonstrates the same two-bit XOR rule. Thus intermediate supervision
changes the learning problem while preserving the required final answer.

Lecture 8 develops three versions of this idea: composing an explicitly
learned XOR gate, exposing the hidden first decision in the path-star task,
and a constructed transformer result in which successive gradient updates
learn retrieval and combination.

\section{Take-home messages}
\label{sec:take-home}

\begin{itemize}[leftmargin=*]
\item Total variation uniformly controls the change in every fixed bounded
downstream cost. A particular cost can remain stable even when this uniform
bound is large.
\item Covariate shift preserves the conditional target; general distribution
shift need not. The distinction depends on the information included in the
input.
\item Under correct specification, a common pointwise optimum survives a
change in input frequencies. Under misspecification, mixture weights determine
the best compromise.
\item Correct inductive bias and coverage address separate requirements. A
density-ratio bound quantifies one sufficient minimum-coverage condition.
\item Teaching and active learning describe different control structures.
The relevant dimensions are what each party knows, who selects each message,
what the communication contains, and whether later messages adapt to earlier
ones.
\item Noiseless examples identify an unknown parity efficiently, and a small
ReLU network represents every fixed parity. For the specified perturbed
population-gradient method, the expected zero--one risk remains within
$e^{-k/36}$ of $1/2$, and the probability of achieving any fixed positive
advantage is at most $1/2+e^{-k/36}$.
\item The lower bound concerns one direct training method. Running-parity
labels preserve the final computation while supplying repeated supervision
for a two-bit rule.
\end{itemize}

\printendnotes

\section*{Bibliographical notes}

The distinction among covariate shift, other forms of dataset shift, and
domain adaptation is developed in the collection edited by
\citet{quinonero2009}. The density-ratio calculation in
\cref{prop:coverage} is elementary; stronger domain-adaptation results replace
the pointwise ratio by discrepancies adapted to a hypothesis class.

The teacher--student terminology follows the separation between active
learning and machine teaching used by \citet{zhu2015} and
\citet{liu2017}. The dimensions in \cref{sec:teacher-student} are the
organizational scheme used in this lecture. Other taxonomies use different
dimensions and terminology.

\Cref{lem:path-comparison,thm:classification} are based on the fixed-parity
lower bound of \citet{shoshani2025}. Lecture 8 gives the Gaussian threshold and trajectory calculations at the
center of their argument. The discussion of
teacher-forced intermediate parity refers to \citet{wen2025}; the path-star
example and the forward-prediction shortcut are from \citet{hu2025}.

\section*{Glossary}

\begin{description}[style=nextline,leftmargin=1.5em,labelindent=0pt]
\item[Active learning]
A protocol in which the student selects queries and receives their labels or
answers from an oracle.
\item[Active teaching]
A protocol in which the teacher adapts its next message to the interaction so
far.
\item[Covariate shift]
A change in the input marginal while the conditional law of the target given
the input remains fixed.
\item[Correct specification]
The model class contains a predictor attaining the optimal conditional
decision at every input in the input space.
\item[Curriculum]
A choice of the order or difficulty of training examples or tasks, fixed in
advance or adapted during training.
\item[Data and task curation]
Selection or weighting of training examples, or modification of the supervised
task, to influence what the student learns.
\item[Domain adaptation]
Learning for a target domain using source data and some information about the
target domain.
\item[Downstream cost]
The expected cost of using a fixed predictor under a specified evaluation or
deployment law.
\item[General distribution shift]
A change in the joint law of input and target, possibly including the
conditional law of the target given the input.
\item[Intermediate supervision]
Training labels reveal intermediate computations whose final result is the
original target behavior.
\item[Inductive bias]
Structure in the model class or learning rule that determines how observations
at some inputs constrain predictions at other inputs.
\item[Machine teaching]
Selection of training information by a teacher that has a target and a model
of the student's learning process.
\item[Minimum coverage]
A condition ensuring that training data identify the degrees of freedom needed
at deployment. \Cref{prop:coverage} uses a stronger, model-independent bounded
density-ratio condition.
\item[Misspecification]
The model class cannot attain the optimal conditional decision at every
input in the input space.
\item[Perturbed gradient descent]
The population-gradient update in \cref{eq:perturbed-gd}, with an independent
Gaussian perturbation added at every step.
\item[Passive learning]
A protocol in which the student receives data selected without using its
queries or current state.
\item[Representation capacity]
The set of input--output functions expressible by a model class, independently
of whether a specified training procedure can find them.
\item[Signed correlation loss]
The objective $\cC(f)=-\E_{P_S}[Yf(X)]$, which isolates the target-dependent
term in the squared-loss and logistic-loss decompositions derived in
\cref{ex:l7-correlation-loss}.
\item[Teacher and student]
The teacher supplies data or feedback. The student updates its state and may
send queries or predictions in a multi-turn protocol.
\item[Total variation]
The uniform discrepancy over event probabilities. Equivalently, it is the
largest possible expectation difference over costs taking values in $[0,1]$.
\item[Zero--one risk]
The probability that a binary-valued predictor differs from the binary label.
\end{description}

\section{Exercises}
\label{sec:exercises}

\exerciseinstructions
\setcounter{courseexercise}{0}
\begin{exercisechunk}
\item \textbf{A common target still requires coverage.}
\label[exercise]{ex:l7-covariate-coverage}
\exerciseprofile{2}{2}{2}{\exproof\quad\exconstruction\quad\exinterpretation}{%
Separate preservation of the pointwise target from identification of that
target by the training law.}
Let \(\cX=\{a,b\}\), \(\cY=\{0,1\}\), and let the loss be zero--one loss.
Suppose the training and deployment laws share the deterministic conditional
kernel
\[
 K(1\mid a)=0,
 \qquad
 K(1\mid b)=1.
\]
The model class contains every function from \(\cX\) to \(\cY\).
\begin{enumerate}[(a),beginpenalty=10000]
\item Prove directly from the two conditional risks that
\(f^\star(a)=0\), \(f^\star(b)=1\) minimizes population risk under every
input marginal on \(\cX\).
\item Let \(\mu_{\rm tr}(a)=1\) and \(\mu_{\rm dep}(b)=1\). Construct a
predictor with zero training risk and deployment risk one. Explain why the
training risk cannot distinguish this predictor from \(f^\star\).
\item Now assume
\(d\mu_{\rm dep}/d\mu_{\rm tr}\le \rho^{-1}\) for some
\(\rho\in(0,1]\). Prove the excess-risk bound in
\cref{prop:coverage} for this finite input space. Identify the step that fails
for the laws in part (b).
\end{enumerate}
\begin{hints}
For (c), write deployment excess risk as a sum over \(x\in\{a,b\}\) and
insert the density ratio into each summand.
\end{hints}
\end{exercisechunk}

\begin{exercisechunk}
\item \textbf{Learning noiseless parity from linear equations.}
\label[exercise]{ex:l7-parity-identification}
\exerciseprofile{2}{2}{2}{\exproof\quad\excalculation}{%
Reduce noiseless parity recovery to linear equations over $\mathbb F_2$ and
prove a high-probability recovery guarantee.}
Let $S\subseteq[d]$ be unknown, and let
$(X^{(1)},Y^{(1)}),\ldots,(X^{(n)},Y^{(n)})$ be independent with law $P_S$
as defined in \cref{sec:parity-setup}.

You may use the following standard linear-algebra fact without proof. Given
$A\in\mathbb F_2^{n\times d}$ and $b\in\mathbb F_2^n$, there is an algorithm
that determines in a polynomial number of arithmetic operations whether
$Az=b$ has a unique solution $z\in\mathbb F_2^d$ and returns that solution
when it is unique. If $z$ is one solution, then it is unique if and only if
\[
Aa\ne0
\qquad\text{for every }a\in\mathbb F_2^d\setminus\{0\}.
\]
The usual linear-system theorem over $\mathbb R$ remains valid with
$\mathbb F_2$ in place of $\mathbb R$ because both are fields.
For background, see Sections 2--4 of David Vogan's
\href{https://math.mit.edu/~dav/gauss.pdf}{notes on Gaussian elimination}.

Construct an algorithm that uses a polynomial number of arithmetic operations
and whose output $\widehat S$ satisfies
\[
\P(\widehat S\ne S)<2^{d-n}.
\]
Then prove that, for every $\delta\in(0,1)$, the algorithm recovers $S$ with
probability at least $1-\delta$ whenever
\[
n\ge d+\left\lceil\log_2\frac1\delta\right\rceil.
\]
\begin{hints}
Write
\[
X_j^{(i)}=(-1)^{U_{ij}},
\qquad
Y^{(i)}=(-1)^{V_i},
\qquad
s_j=\I\{j\in S\}.
\]
With arithmetic over $\mathbb F_2$, the examples give
$\mathbf v=\mathbf U\mathbf s$. Apply the stated linear-system procedure. To
bound its failure probability, first fix a nonzero
$a\in\mathbb F_2^d$ and calculate $\P(\mathbf Ua=0)$. Then use a union bound
over the $2^d-1$ possible choices of $a$.
\end{hints}
\end{exercisechunk}

\begin{exercisechunk}
\item \textbf{Representing parity with one ReLU layer.}
\label[exercise]{ex:l7-parity-representation}
\exerciseprofile{2}{2}{2}{\exproof\quad\excalculation}{%
Verify that a width-$(|S|+1)$ ReLU network represents a fixed parity exactly.}
Fix $S\subseteq[d]$, write $k=|S|$, and let $\mathbf 1_S\in\{0,1\}^d$ be
the indicator vector of $S$. Consider
\begin{equation*}
N(x)=[-\mathbf 1_S^\top x+k+1]_+
+\sum_{j=1}^k4j(-1)^j
[-\mathbf 1_S^\top x+k+1-2j]_+.
\end{equation*}
Prove that
\[
N(x)=\chi_S(x)
\qquad\text{for every }x\in\{-1,1\}^d.
\]
Conclude that the network in \cref{prop:parity-representation} needs at most
$k+1$ hidden units.

\begin{hints}
Let $a$ be the number of negative coordinates of $x$ indexed by $S$. First
show that $-\mathbf 1_S^\top x+k+1=2a+1$ and identify which ReLU terms are
active. If $\widetilde N(a)$ denotes the resulting value, compare
$\widetilde N(a)$ with $\widetilde N(a-1)$. Pair consecutive terms when
evaluating $\sum_{j=1}^{a-1}j(-1)^j$.
\end{hints}
\end{exercisechunk}

\begin{exercisechunk}
\item \textbf{What signed correlation controls.}
\label[exercise]{ex:l7-correlation-loss}
\exerciseprofile{2}{2}{2}{\exproof\quad\excalculation\quad\exconstruction\quad\exinterpretation}{%
Relate signed correlation to classification and to the gradients of standard
surrogate losses.}
Under the fixed-parity law $P_S$, recall
\[
\cC(f)=-\E_{P_S}[Yf(X)].
\]
\begin{enumerate}[(a),beginpenalty=10000]
\item Let $h:\{-1,1\}^d\to\{-1,1\}$. Prove
\[
\cC(h)=2R_{01}(h)-1,
\qquad
R_{01}(h)=P_S(h(X)\ne Y).
\]
\item For $\varepsilon>0$, set $f_\varepsilon=\varepsilon\chi_S$. Compute
$\cC(f_\varepsilon)$ and $R_{01}(\operatorname{sign}(f_\varepsilon))$.
Explain why a correlation value close to zero does not determine the
classification risk of the sign of a real-valued score.
\item For any score $f:\{-1,1\}^d\to\R$, derive
\[
\E_{P_S}[(f(X)-Y)^2]
=1+\E_{P_S}[f(X)^2]+2\cC(f).
\]
\item Derive
\[
\E_{P_S}[\log(1+e^{-Yf(X)})]
=\E_{P_S}\!\left[\log\!\left(2\cosh\frac{f(X)}2\right)\right]
+\frac12\cC(f).
\]
If $f=f_\theta$ is differentiable in $\theta$, differentiate the identities
in parts (c) and (d) and identify every term in addition to
$\nabla_\theta\cC(f_\theta)$.
\end{enumerate}
\begin{hints}
For (a), determine the value of $Yh(X)$ on the correct and incorrect events.
For (d), verify the pointwise identity
\[
\log(1+e^{-z})=\log\!\left(2\cosh\frac z2\right)-\frac z2.
\]
\end{hints}
\end{exercisechunk}


\clearpage
\begin{thebibliography}{99}

\bibitem[Hu et al.(2025)]{hu2025}
Edward S. Hu, Kwangjun Ahn, Qinghua Liu, Haoran Xu, Manan Tomar,
Ada Langford, Dinesh Jayaraman, Alex Lamb, and John Langford.
\newblock Learning to achieve goals with belief state transformers.
\newblock 2025. \url{https://arxiv.org/abs/2410.23506}.

\bibitem[Wen et al.(2025)]{wen2025}
Kaiyue Wen, Huaqing Zhang, Hongzhou Lin, and Jingzhao Zhang.
\newblock From sparse dependence to sparse attention: Unveiling how
chain-of-thought enhances transformer sample efficiency.
\newblock International Conference on Learning Representations, 2025.
\newblock \url{https://arxiv.org/abs/2410.05459}.

\bibitem[Liu et al.(2017)]{liu2017}
Weiyang Liu, Bo Dai, Ahmad Humayun, Charlene Tay, Chen Yu, Linda B. Smith,
James M. Rehg, and Le Song.
\newblock Iterative machine teaching.
\newblock International Conference on Machine Learning, 2017.
\newblock \url{https://proceedings.mlr.press/v70/liu17b.html}.

\bibitem[Qui\~nonero-Candela et al.(2009)]{quinonero2009}
Joaquin Qui\~nonero-Candela, Masashi Sugiyama, Anton Schwaighofer, and
Neil D. Lawrence, editors.
\newblock \emph{Dataset Shift in Machine Learning}.
\newblock MIT Press, 2009.

\bibitem[Shoshani and Shamir(2025)]{shoshani2025}
Itamar Shoshani and Ohad Shamir.
\newblock Hardness of learning fixed parities with neural networks.
\newblock 2025. \url{https://arxiv.org/abs/2501.00817}.

\bibitem[Zhu(2015)]{zhu2015}
Xiaojin Zhu.
\newblock Machine teaching: An inverse problem to machine learning and an
approach toward optimal education.
\newblock AAAI Conference on Artificial Intelligence, 2015.

\end{thebibliography}

\end{document}
